% ******************************************************** %
%              TEMPLATE DE INFORME ORGA2 v0.1              %
% ******************************************************** %
% ******************************************************** %
%                                                          %
% ALGUNOS PAQUETES REQUERIDOS (EN UBUNTU):                 %
% ========================================
%                                                          %
% texlive-latex-base                                       %
% texlive-latex-recommended                                %
% texlive-fonts-recommended                                %
% texlive-latex-extra?                                     %
% texlive-lang-spanish (en ubuntu 13.10)                   %
% ******************************************************** %


\documentclass[a4paper]{article}
\usepackage[spanish]{babel}
\usepackage[utf8]{inputenc}
\usepackage{charter}   % tipografia
\usepackage{tipa}
\usepackage{graphicx}
%\usepackage{makeidx}
\usepackage{paralist} %itemize inline

%\usepackage{float}
\usepackage{amsmath, amsthm, amssymb}
%\usepackage{amsfonts}
%\usepackage{sectsty}
%\usepackage{charter}
%\usepackage{wrapfig}
%\usepackage{listings}
%\lstset{language=C}

% \setcounter{secnumdepth}{2}
\usepackage{underscore}
\usepackage{caratula}
\usepackage{url}


% ********************************************************* %
% ~~~~~~~~              Code snippets             ~~~~~~~~~ %
% ********************************************************* %

\usepackage{color} % para snipets de codigo coloreados
\usepackage{fancybox}  % para el sbox de los snipets de codigo

\definecolor{litegrey}{gray}{0.94}

\newenvironment{codesnippet}{%
	\begin{Sbox}\begin{minipage}{\textwidth}\sffamily\small}%
	{\end{minipage}\end{Sbox}%
		\begin{center}%
		\vspace{-0.4cm}\colorbox{litegrey}{\TheSbox}\end{center}\vspace{0.3cm}}



% ********************************************************* %
% ~~~~~~~~         Formato de las páginas         ~~~~~~~~~ %
% ********************************************************* %

\usepackage{fancyhdr}
\pagestyle{fancy}

%\renewcommand{\chaptermark}[1]{\markboth{#1}{}}
\renewcommand{\sectionmark}[1]{\markright{\thesection\ - #1}}

\fancyhf{}

\fancyhead[LO]{Sección \rightmark} % \thesection\ 
\fancyfoot[LO]{\small{Nombre Apellido, Nombre Apellido, Nombre Apellido}}
\fancyfoot[RO]{\thepage}
\renewcommand{\headrulewidth}{0.5pt}
\renewcommand{\footrulewidth}{0.5pt}
\setlength{\hoffset}{-0.8in}
\setlength{\textwidth}{16cm}
%\setlength{\hoffset}{-1.1cm}
%\setlength{\textwidth}{16cm}
\setlength{\headsep}{0.5cm}
\setlength{\textheight}{25cm}
\setlength{\voffset}{-0.7in}
\setlength{\headwidth}{\textwidth}
\setlength{\headheight}{13.1pt}

\renewcommand{\baselinestretch}{1.1}  % line spacing

% ******************************************************** %


\begin{document}


\thispagestyle{empty}
\materia{Organización del Computador II}
\submateria{Primer Cuatrimestre de 2019}
\titulo{Trabajo Práctico II}
\subtitulo{subtítulo del trabajo}
\integrante{Nombre}{XXX/XX}{mail}
\integrante{Nombre}{XXX/XX}{mail}

\maketitle
\newpage

\thispagestyle{empty}
\vfill
\begin{abstract}
En el presente trabajo se describe la problemática de ...
\end{abstract}

\thispagestyle{empty}
\vspace{3cm}
\tableofcontents
\newpage


%\normalsize
\newpage

\section{Introducción}

SIMD (single instruction, multiple data) es una arquitectura de computaci\'on centrada en el procesamiento simultaneo de una operaci\'on sobre varios datos mediante una \'unica isntrucci\'on. Tal t\'ecnica suele ser aplicada para trabajar de manera eficiente sobre vectores y matrices de datos. \par
El objetivo de este trabajo pr\'actico es consolidar los conocimientos adquiridos en las clases sobre SIMD. Utilizando esta arquitectura implementaremos distintos algoritmos en ASM creando filtros sobre imágenes para luego comparar su rendimiento contra su implementación en C.


\section{Desarrrollo}

\begin{figure}
  \begin{center}
	\includegraphics[scale=0.66]{img/logouba.jpg}
	\label{nombreparareferenciar}
  \end{center}
\end{figure}

\subsection{Imagen Fantasma}

\subsubsection{Descripci\'on}

\paragraph{} Este filtro consiste en generar una imagen fantasma sobre la original, para ello se utiliza la misma imagen original al doble de tamaño y en escala de grises. Ademas se generara a partir de un punto cuyas coordenadas nos dan por parametro.

\subsubsection{Implementaci\'on}

\paragraph{} En esta implementaci\'on se realizan dos recorridos diferentes por cada iteraci\'on sobre la matriz $src$: 
\begin{itemize}
	\item Uno de los recorridos es similar al implementado en C, es decir recorriendo todos los pixeles de cada fila. Con la principal diferencia de que se levantan 4 pixeles en vez de 1.
	\item En nuestra implementaci\'on no se calcula $ii$ ni $jj$, en cambio existe un segundo recorrido que levanta los pixeles correspondientes a estos indices.
\end{itemize}

\paragraph{Recorrido de $ii$ y $jj$:} 
\begin{center}
	Sabemos que $ii = i/2 + offset_x$ , \, $jj = j/2 + offset_y$ representando $"/"$ la divisi\'on entera. Fijando el $offset_y$ en 0 veamos como cambia $jj$ a partir de una sucesi\'on de valores para $j$:
	
	\begin{align*}
		j &= [0, 1, 2, 3, 4, 5, 6, 7, 8, 9]\\
		jj &= [0, 0, 1, 1, 2, 2, 3, 3, 4, 4]
	\end{align*}    
	
	\paragraph{} El valor de $jj$ para los $j$ impares es el mismo que para el par menor y mas cercano a ese $j$ impar, ya que el operador $"/"$ descarta el resto y devuelve solo el cociente es decir: $j = 2*q \, \wedge \, j+1 = 2*q + r \Rightarrow j/2 = q$ \, y \, $(j+1)/2 = q$. Podemos usar esta informaci\'on para afirmar que si $j = 2 * q \Rightarrow jj(j) == jj(j+1)$ (notar que en el calculo de $jj$, $offset_y$ es siempre la misma constante sumada, no cambia la igualdad).
	
	\paragraph{} Esto lo aprovechamos para no calcular $ii$ ni $jj$ en cada iteraci\'on, ya que $i, j$ empiezan en $0 \Rightarrow ii = offset_x, jj=offset_y$. Luego basta con incrementar en 1 a $ii$ o $jj$ por cada 2 incrementos de $i$ y $j$ respectivamente (en nuestro caso los incrementos son de a 2 para $ii$, $jj$ y 4 para $i$ y $j$ ya que se procesan 4 pixeles).
\end{center}


\paragraph{}Para calcular $B$, utilizamos una mascara que re ordena las componentes de los 2 pixeles $m[ii][jj]$ y $m[ii][jj+1]$ de la siguiente manera: $blue, green, rojo, green$. Este nuevo orden permite que realicemos una suma horizontal, sin embargo no existe suma horizontal de a Byte, por lo que adem\'as extendemos con 0 cada byte usando la misma mascara.  

\begin{align*}
	xmm &= [\, \, \, \, * \, \, \, \, \, ,\, \, \, \, \, * \, \, \, \, \, ,\, \, \, \, * \, \, \, \, \, ,\, \, \, \, * \, \, \, \, \, ,\, \, \, \, * \, \, \, \, \, \, ,\, \, \, \, * \, \, \, \, \, ,\, \, \, \, * \, \, \, \, \, ,\, \, \, \, * \, \, \, \, \, \, , \, \, \, \,  A_1 \, , \, \, \, \, R_1 \, , \, \, \, \, G_1 \, , \, \, \, \, B_1 \, , \, \, \, \, A_0 \, , \, \, \, \, R_0 \, , \, \, \, \, G_0 \, , \, \, \, \, B_0 \, \,]\\
	mascara &= [0x80, 0x05, 0x80, 0x06, 0x80, 0x05, 0x80, 0x04, 0x80, 0x01, 0x80, 0x02, 0x80, 0x01, 0x80, 0x00] \\
\end{align*}
	\hrule
\begin{align*}
	pshufb &= [G_1, R_1, G_1, B_1, G_0, R_0, G_0, B_0] \leftarrow \text{en words}
\end{align*} 

\paragraph{}Una vez ordenadas las componentes se ejecutan 2 sumas horizontales para generar $B_0$ y $B_1$ en los dos words menos significativos del registro. Separamos los $B$ en registros diferentes ($R_{b0}$, $R_{b1}$) y desempaquetamos a doublewords cada uno, adem\'as replicamos cada $B$ en los 4 doubleword en su respectivo registro. Luego los convertimos en floats para dividirlos por 8 (b/4/2 == b/8).

\paragraph{} Para procesar los pixeles, realizamos los siguiente pasos:

\begin{itemize}
	\item Se mueven los 4 pixeles a 4 registros difrentes: $R_{px0}$, $R_{px1}$, $R_{px2}$ y $R_{px3}$.
	\item En cada registro se realizan las conversiones necesarias para pasar de bytes a words y de words a doublewords para el pixel que le corresponda.
	\item Convertimos las componentes de cada registro a floats.
	\item Realizamos multiplicaciones vectoriales entre una constante con 4 datos de tipo float con valor 0,9 y cada uno de los $R_{px}$.
	\item Sumamos $R_{px0}$, $R_{px1}$ con $R_{b0}$ y $R_{px2}$, $R_{px3}$ con $R_{b1}$.
\end{itemize}

Con esto tenemos nuestros 4 pixeles procesados. Por \'ultimo queda realizar 4 conversiones de float con truncamiento a doubleword en $R_{px0}$, $R_{px1}$, $R_{px2}$, desempaquetar (haciendo las conversiones intermedias) cada registro en 1 solo que junte los 4 pixeles procesados.

\paragraph{Comentarios sobre el signado en operaciones:} Algunas conversiones u operaciones no existen para n\'umeros sin signo, esto puede generar problemas si lo n\'umeros sin signo sobre los que se quiere operar son tan grandes que su representaci\'on binaria tiene el bit mas significativo en 1, pues en su interpretaci\'on signada serian negativos.
\paragraph{}Los n\'umeros mas grande sobre los que trabajamos son $B$ y $componente(px) * 0,9 + B/2$. Si todas las componentes de $m[ii][jj]$ valen 255 $\Rightarrow B = 4*255/4 = 255$. Luego $componente(px) * 0,9 + B/2 \simeq  230 + 126 \simeq 356$. Como todas las operaciones son realizadas luego de convertir los datos al menos a words tenemos los suficientes bits para representar $356$ sin usar el bit mas significativo.

\subsection{Color Bordes}
\subsubsection{Descripci\'on}

\paragraph{} Color bordes modifica cada pixel operando mediante restas sobre los 8 pixeles que lo rodean, generando valores pequeños y similares para cada componente si los 8 pixeles tenían un color parecido entre ellos y variando estos valores en caso contrario. Esto genera diferentes componentes conexas con colores oscuros separadas por bordes de colores variables y claros. Adem\'as agrega un marco blanco de 1 pixel.  

\subsubsection{Desarrollo}

\paragraph{} La matriz se recorre de manera similar a la implementaci\'on en C, pero en cada iteraci\'on se modifican 2 pixeles. Utilizamos las herramientas de SIMD para levantar 3 filas ($i-1$, $i$ y $i+1$) de 4 pixeles, re ordenar los datos y luego operar. Nos referiremos a cada una de estas filas con mediante los nombres $R_{i-1}$, $R_i$ y $R_{i+1}$. Llamaremos $px0$, $px1$, $px2$ y $px3$ a los pixeles de cada fila. Modificaremos $px1$ y $px2$ de $R_i$.

\paragraph{ii:} Con el objetivo de realizar $m[ii][j-1] - m[ii][j+1]$ mediante restas horizontales armamos 2 mascaras, una para procesar los datos de $px1$ y otra para $px2$. Esta mascara adem\'as extiende los datos de byte a words con ceros.
 


\begin{align*}
	R_i &= [\, \, \, \, A_3 \, ,\, \, \, \, \, R_3 ,\, \, \, \, G_3 \, ,\, \, \, \, B_3 \, ,\, \, \, \, A_2 \, ,\, \, \, \, R_2 \, ,\, \, \, \, G_2 \, ,\, \, \, \, B_2 \, , \, \, \, \,  A_1 \, \, , \, \, \, \, R_1 \, , \, \, \, \, G_1 \, , \, \, \, \, B_1 \, , \, \, \, \, A_0 \, \, , \, \, \, \, R_0 \, , \, \, \, \, G_0 \, , \, \, \, \, B_0 \, \,]\\
	indice &= [0x0f, 0x0e, 0x0d, 0x0c, 0x0b, 0x0a, 0x09, 0x08, 0x07, 0x06, 0x05, 0x04, 0x03, 0x02, 0x01, 0x00]\\
	mascara_{px1} &= [0x80, 0x05, 0x80, 0x06, 0x80, 0x05, 0x80, 0x04, 0x80, 0x01, 0x80, 0x02, 0x80, 0x01, 0x80, 0x00] \\
\end{align*}
\hrule
\begin{align*}
	pshufb &= [0, 0, R_2, R_0, G_2, G_0, B_2, B_0] \leftarrow \text{en words}
\end{align*} 

\paragraph{} El proceso anterior es aplicado a cada uno de los 3 registros con ambas mascaras. Teniendo los datos acomodados en 6 registros (3 filas para cada pixel) los siguientes pasos son directos y simples. Primero realizamos una resta horizontal en cada uno de ellos consiguiendo $m[ii][j-1] - m[ii][j+1]$ de cada componente, luego aplicamos el valor absoluto y acumulamos los valores de cada componente.

\paragraph{jj :}En este caso primero operamos y luego reordenamos los datos. Esta vez nuestro iterador es $jj$ y se mueve por $j$ mientras que opera entre las filas $i-1$ e $i+1$. Queremos calcular $m[i-1][jj] - m[i+1][jj]$ para cada componente.

\begin{itemize}
	\item Desempaquetamos $R_{i-1}$ y $R_{i+1}$ a words.
	\item Restamos $R_{i-1}$ y $R_{i+1}$ desempaquetados.
	\item Aplicamos valor absoluto a los resultados
	\item A partir de este punto tenemos 2 registros $R_1$ y $R_2$. $R_1$ contiene las restas en words de los componentes de $j$ y $j-1$, mientras $R_2$ contiene las restas de $j+2$ y $j+1$.
	\item Para acumular estos resultados utilizo 1 registro intermedio para cada uno. Paso $R_1$ a $R_3$ y shifteo los resultados $j$ a la parte mas baja de $R_3$. Ahora el primer word de $R_1$, $R_2$ y $R_3$ contienen los resultados de $j-1$, $j+1$ y $j$ respectivamente.
	\item Sumamos los 3 registros y obtenemos las acumulaciones $jj$ de las operaciones entre componentes para $px1$. Haciendo un proceso similar, pero con shifteos a la parte alta obtenemos las acumulaciones $jj$ de $px2$. 
\end{itemize}

\paragraph{} Finalmente sumamos las acumulaciones $ii$ y $jj$ para $px1$ y $px2$, convertimos los resultados utilizando empaquetado saturado sin signo y lo guardamos en la imagen destino.


\subsection{ReforzarBrillo}

\subsubsection{Descripción}

Este filtro aumenta y disminuye el brillo de una imagen según el brillo que ya tenia. Sí el
pixel supera un umbral, el brillo se aumenta, si el pixel esta debajo de un umbral, el brillo se
disminuye. El efecto resultante es un refuerzo del brillo diferenciado.

\subsubsection{Implementaci\'on}

En este filtro consideramos varias cosas a la hora de implementar este filtro.

Para poder detenerminar si un pixel superaba un estaba por debajo de un umbral, necesitabamos hacer la siguiente cuenta:

\begin{center}
	b = ( pixel.R + 2 * pixel.G + pixel.B ) / 4
\end{center}

Seguido a esto, ya que cada componente de un pixel mide 1 byte, dando que el tama\~no total de un pixel es de 4 bytes, en un registro xmm podemos levantar 4 pixeles. Sin embargo, si queremos que nuestra cuenta sea exacta, fue necesario extender el tama\~no de cada componente a 2 bytes. Para esto, en vez de utilizar una instruccion de extension de como pmovzxbw, decidimos utilizar la instrucción de pshufb con una mascara que nos dejá de la componente Verde en la posicion en la que deberia estar la transparencia. Quedando el registro xmm de la siguiente forma: 

\begin{center}
	xmm = [ G1 \textpipe\ R1 \textpipe\ G1 \textpipe\ B1 \textpipe\ G0 \textpipe\ R0 \textpipe\ G0 \textpipe\ B0 ]
\end{center}

Esto nos facilit\'o la suma, ya que solo tuvimos que usar la suma horizontal quedando el registro xmm:

\begin{center}
	xmm = [ 0 \textpipe\ 0 \textpipe\ 0 \textpipe\ 0 \textpipe\ G1 + R1 \textpipe\ G1 + B1 \textpipe\ G0 + R0 \textpipe\ G0 + B0 ]
\end{center}

y si usamos nuevamente la suma horizontal nuevamente obtendremos:

\begin{center}
	xmm = [ 0 \textpipe\ 0 \textpipe\ 0 \textpipe\ 0 \textpipe\ 0 \textpipe\ 0 \textpipe\ R1 + 2 $\times$ G1 + B1 \textpipe\ R0 + 2 $\times$ G0 + B0 ]
\end{center}

Finalmente divisimos por 4 usando un shift empaquetado de a word de 2 bits a derecha para que nos quede b, como nosotros queriamos.

Ahora, para poder compararlos con los umbrales que pueden ser numeros de hasta 4 bytes, es necesario extender nuevamente de word a dword (2 a 4 bytes) quedando los b de cada pixel expresados en dwords.

Una vez que tenemos nuestros b, pasamos a analizar cual b es mayor o menor a los umbrales pasados por parametro. Sabemos que como estamos preguntando por mayor o menor estrictos, no es posible que un b cumpla ambas condiciones al mismo tiempo (?????? estoy super exceptico con respecto a esto).

\section{Resultados}
\subsection{Planteo de hipótesis y experimentos}

\subsubsection{Comparaci\'ones con C}

Con el objetivo saber cual de las implementaciones es mas r\'apida compararemos los filtros hechos en ASM utilizando SIMD con las implementaciones en C con diferentes niveles de optimizaciones. Creemos que ASM  sera mucho mas rapido, incluso con optimizaciones aplicadas al codigo C. re escribir mas lindo (?) 

\subsubsection{Desenrollado de ciclos}

\subsubsection{Accesos a memoria}



\subsection{Dise\~no de instancias y aspectos t\'ecnicos}


\paragraph{} Para los experimentos de comparaci\'on de rendimiento armamos instancias de imagenes cuyas resoluciones son multiplo de $32*i x 18*i$ con $\forall i \in [1..100]$.


La computadora en donde corrimos los experimentos cuenta con un procesador Intel Core I5-6200U, con 3MB de cach\'e y $2,30$GHz de frecuencia.


Para evitar outliers corrimos de 500 veces cada imagen en todos los experimentos. Adem\'as en el caso de los experimentos de desenrollado de ciclos y costo de accesos a memoria decidimos analizar las diferencias entre las medianas de la resoluci\'on de un filtro original y su modificaci\'on para tener una mejor visualizaci\'on de los datos.


\section{Conclusión}


\end{document}

